\section{Scalar convex geometry at finite accuracy}\label{sec:scalar-geometry}

For one convex output on an interval, chord geometry gives a uniform
comparison with arbitrary convex integer lifts: two extra binaries
suffice at every positive accuracy. We first prove this geometric
statement, then develop the certified computations needed to obtain a
short rational formulation. The latter requirement is substantive:
a partition with a small logarithm of its cell count may still be too
long to enumerate.

For a continuous convex function $f:[0,1]\to\R$, let
$N_\eta(f)$ be the smallest number of intervals in a partition of
$[0,1]$ on each of which the chord of $f$ has error at most $\eta>0$.
This number is finite by uniform continuity. Chord errors decrease when
an interval is restricted: the chord on a subinterval lies below the
original chord and above $f$. Affine terms change neither these errors
nor the formulation dimensions.

\begin{lemma}[Chord refinement and scalar packing]\label{lem:scalar-chords}
For every continuous convex $f$ and $\eta>0$,
\begin{equation}\label{eq:scalar-chord-counts}
 N_\eta(f)\le3N_{2\eta}(f),\qquad
 N_{2\eta}(f)\le2^{\pconv(f;\eta)},\qquad
 \pbin(f;\eta)\le\lceil\log_2N_\eta(f)\rceil.
\end{equation}
Consequently
\begin{equation}\label{eq:scalar-two-bit}
 \pconv(f;\eta)\le\pbin(f;\eta)\le\pconv(f;\eta)+2.
\end{equation}
There is a finite maximal set of $P\ge1$ points pairwise satisfying
$J_f(a,b)>\eta$. Every such set obeys
\begin{equation}\label{eq:scalar-packing-chords}
 N_{2\eta}(f)\le2P,\qquad N_\eta(f)\le6P.
\end{equation}
These are existence statements with real coefficients and unrestricted
continuous size.
\end{lemma}
\begin{proof}
On an interval of chord error at most $2\eta$, write $g$ for the chord
minus $f$. It is continuous, concave, nonnegative, and zero at both
endpoints. If $\max g>\eta$, its superlevel set $\{g\ge\eta\}$ is an
interior interval $[u,v]$, with $g(u)=g(v)=\eta$. Split at $u,v$.
On the outside pieces the new error, $g$ minus its own chord, is at
most $g\le\eta$; on the middle piece it equals $g-\eta\le\eta$.
If $\max g\le\eta$ no split is needed. This proves the first inequality.

A parity contact from \cref{lem:parity} has compact span $[a,b]$.
The concave chord gap on this span has maximum at most twice its
midpoint value: concavity between a maximizer and the farther endpoint
proves this bound. Thus its chord error is at most $2\eta$.
The at most $2^p$ spans cover $[0,1]$. A finite interval cover can be
trimmed to a partition with no more intervals: starting at zero, take
a covering interval reaching farthest to the right and continue;
each resulting piece is contained in a covering member. Hence
$N_{2\eta}\le2^p$. Each chord band $s-\eta\le w\le s$ contains
the graph and has error at most $\eta$. These bounded polyhedra have
the binary encoding of \cref{lem:disjunction}, proving the third
inequality and \eqref{eq:scalar-two-bit}.

Uniform continuity bounds the cardinality of every incompatible set,
since sufficiently close points have gap at most $\eta$. Add points
until a finite maximal set is reached. For fixed $x$, its compatible
set $\{y:J_f(x,y)\le\eta\}$ is a closed interval containing $x$.
Indeed $J_f(a,b)$ increases under enlargement of $[a,b]$, by monotone
secant slopes; equivalently, smooth approximation reduces this to
$\partial_aJ_f\le0$ and $\partial_bJ_f\ge0$. Maximality says these
$P$ intervals cover $[0,1]$. Splitting each at its selected point gives
at most $2P$ intervals of endpoint midpoint gap at most $\eta$ and
chord error at most $2\eta$. Trim the cover and apply three-piece
refinement to obtain \eqref{eq:scalar-packing-chords}.
\end{proof}

The factor-two chord/midpoint comparison is also explicit in
\citet[Lemma 2.1]{simchowitz2018}. Their accuracy-dependent local
approximation modulus and endpoint corrections concern sampling
complexity; the dimension comparison here uses convex integer lifts.
The comparison uses the established parity and disjunction mechanisms
\citep{lubin2022,vielma2018}. Efficiently computing an exponentially
long partition is a separate requirement; the finite two-bit bound
alone supplies no polynomial-size rational formulation.

\subsection{An accuracy-dependent curvature measure}

Suppose $f\in C^2([0,1])$ and $g=f''\ge0$ is nondecreasing. Define
\begin{equation}\label{eq:curvature-density}
 A_\eta(x)=\sqrt{g(x)/\eta},\quad
 \rho_\eta(x)=\min\{A_\eta(x),(1-x)A_\eta(x)^2\},\quad
 M_\eta=\int_0^1\rho_\eta(x)\,dx.
\end{equation}
The linear branch suppresses curvature contributions too small to
require separate intervals at the specified accuracy.

\begin{theorem}[Truncated curvature mass]\label{thm:curvature-mass}
Under these assumptions,
\begin{equation}\label{eq:mass-chord-law}
 N_\eta\le2M_\eta+1,\qquad M_\eta\le24N_\eta,
\end{equation}
and
\begin{equation}\label{eq:mass-dimension-law}
 \log_2(1+M_\eta)-\log_2 49\le\pconv(f;\eta)
 \le\pbin(f;\eta)\le\log_2(1+M_\eta)+2.
\end{equation}
For every $[a,b]\subset[0,1]$, with
$m=\int_a^b\rho_\eta$ and
$E=\int_a^b(b-t)A_\eta(t)^2\,dt$, one has
\begin{equation}\label{eq:mass-local-remainder}
 E\le m^2+\tfrac32m.
\end{equation}
\end{theorem}
\begin{proof}
Suppress the subscript on $A$. On the set $(b-t)A(t)<1$, the
integrand of $E$ is at most $\rho_\eta(t)$. On its complement, put
$z=(b-t)A(t)\ge1$. Monotonicity gives
\[
 m\ge\int_t^b\min\{A(t),(b-s)A(t)^2\}\,ds=z-\tfrac12.
\]
There $\rho_\eta(t)=A(t)$, so the contribution to $E$ is at most
$(m+1/2)m$. This proves \eqref{eq:mass-local-remainder}. If $m\le1/2$,
then $E\le1$. Convexity bounds the chord error by the left-end Taylor
remainder $\eta E$. Dividing total mass into pieces of mass at most
$1/2$, and absorbing zero-density intervals, proves the first bound
in \eqref{eq:mass-chord-law}. Zero mass means $f$ is affine.

For the reverse estimate, take an admissible interval, put
$\ell=b-a$, $c=(a+b)/2$, and write $J=J_f(a,b)\le\eta$.
The right half of the tent-kernel representation gives
$J\ge g(c)\ell^2/16$. The left half of the Taylor remainder is at
most $3g(c)\ell^2/8\le6J$, and its right half is at most $2J$.
Thus $E\le8$. Introduce the nonnegative potential
$\chi(x)=(1-x)A(x)$. We claim for every interval that
\begin{equation}\label{eq:mass-potential}
 m\le\chi(b)-\chi(a)+2E+2\sqrt{2E}.
\end{equation}
Put $d=1-b$. If $\ell\le d$, then $m\le\ell A(b)$ and
\[
 m-\chi(b)+\chi(a)
 \le2\ell A(a)-(d-\ell)(A(b)-A(a))
 \le2\ell A(a)\le2\sqrt{2E}.
\]
If $\ell>d$, split at $b-d$. On the left part,
$\rho_\eta(t)\le(1-t)A(t)^2\le2(b-t)A(t)^2$, whose integral is at
most $2E$. On the right its integral is at most $dA(b)=\chi(b)$.
Hence $m-\chi(b)+\chi(a)\le2E+\chi(a)
\le2E+2\ell A(a)\le2E+2\sqrt{2E}$. This includes $b=1$.
On an admissible interval \eqref{eq:mass-potential} is at most the
potential difference plus $24$. Summing over a partition telescopes,
with $\chi(1)=0\le\chi(0)$, and proves $M_\eta\le24N_\eta$.

Pointwise $\rho_\eta\le2\rho_{2\eta}$. By
\cref{lem:scalar-chords}, $M_\eta\le48\,2^{\pconv(f;\eta)}$.
Since $p\ge0$, this yields the lower dimension bound. The upper follows
from $N_\eta\le2(1+M_\eta)$ and binary chord-band encoding.
\end{proof}

\begin{example}[Curvature and allocation obstructions]
\label{ex:positive-degree-gap}
Let $f_M(x)=M^{-1}\sum_{j=1}^M x^{2^j}$, of degree $D=2^M$.
At fixed $\eta=1/16$,
\begin{equation}\label{eq:raw-curvature-obstruction}
 \int_0^1\sqrt{f_M''(x)}\,dx\ge\frac{\sqrt M}{4\sqrt2},
 \qquad \pconv\le\pbin\le4.
\end{equation}
Indeed, on the disjoint layer
$[1-1/k,1-1/(2k)]$, $k=2^j$, the $j$th monomial contributes curvature
at least $k^2/(8M)$. Its length is $1/(2k)$; summing the resulting
integrals proves the lower bound. The curvature estimate uses
$k(k-1)\ge k^2/2$ and $(1-1/k)^{k-2}\ge1/4$.
The function increases from zero to one; inverse images of the sixteen
equal output intervals yield sixteen chords with error at most $1/16$.
Thus raw curvature arclength cannot characterize finite-accuracy count
uniformly over functions. In contrast,
$M_\eta\le\eta^{-1}[f_M(1)-f_M(0)-f_M'(0)]=\eta^{-1}$.

At the different tolerance $\eta=1/(512M)$, the same family satisfies
\begin{equation}\label{eq:degree-allocation-obstruction}
 \lceil\log_2M\rceil\le\pconv\le\pbin
 \le\lceil\log_2(512M)\rceil.
\end{equation}
For the lower bound, choose $x_j=1-2^{-j}$. If $j<\ell$ and $k=2^j$,
then $x_\ell-x_j\ge1/(2k)$ and the selected monomial has curvature
at least $k^2/8$ throughout that interval. Its midpoint gap is at least
$1/256$, so $J_{f_M}(x_j,x_\ell)\ge1/(256M)>\eta$. Parity proves
the claim. The upper again partitions the range, now into $512M$ pieces.
The coefficient-sum allocation defined in \cref{sec:positive-allocation}
has $\Phi=\frac12\log_2(512M)$ here. Consequently
$\pconv-\Phi=\frac12\log_2\log_2D+O(1)$, with matching upper
and lower orders. This limits that allocation benchmark; it does not
force a degree penalty relative to the true minimum.
\end{example}

\section{Certified scalar compilation}\label{sec:scalar-compilation}

A short indexed computation can describe exponentially many intervals.
The construction has three parts: compute one cell from its index,
interpolate its endpoints with continuous variables, and add a certified
vertical band containing the graph. Only the cell index is declared
integral. The geometric lower bounds then compare this count with all
convex integer lifts, including those with unrestricted continuous size.
The following established circuit principle specifies exactly which
variables count as integers. Circuit-to-polyhedron compilation is explicit
in \citet[Section 3, Lemma 1]{avis2019}; products of discrete function
values with continuous weights also have direct base-two formulations
\citep[Section 2]{adams2012}.

\begin{lemma}[Indexed endpoint compiler]\label{lem:indexed-compiler}
Fix rational instance data and an $L$-bit index, possibly restricted to
$0\le k<K\le2^L$. Suppose a deterministic algorithm computes the two
endpoint vectors of cell $k$ in time polynomial in the instance length
and $L$, with polynomial output length and a known common denominator
of polynomial bit length for each output coordinate. Their line segment, with one common
interpolation parameter, has a polynomial-size rational linear
formulation using only the $L$ index binaries. All internal wires and
interpolation products are continuous variables. The construction is
uniform in the stated running-time bound.
\end{lemma}
\begin{proof}
Unroll the bounded computation into a polynomial-size Boolean circuit,
with the instance as constants. Fix its worst-case running time, padding
early exits with flags. Use continuous wires in $[0,1]$, NOT equations
$v=1-u$, and AND inequalities
\[
 0\le v,\quad v\le u_1,\quad v\le u_2,\quad v\ge u_1+u_2-1.
\]
Induction along the acyclic circuit forces every wire to its exact
Boolean value for integral input bits. A comparison circuit excludes
invalid indices. For $0\le\theta\le1$ and every endpoint output bit
$v$, impose $0\le t\le\theta$, $t\le v$, $t\ge\theta+v-1$.
Then $t=\theta v$ exactly. Decode the fixed-denominator endpoint
numerators linearly and form $(1-\theta)a_k+\theta b_k$ using those
products. Signed outputs are represented by a fixed rational offset
and a nonnegative numerator. A fixed-width rational denominator can
be incorporated in the linear decoding coefficients. No gate requires
an additional integrality declaration.
\end{proof}

\subsection{Certified integration of monotone polynomial curvature}

The next proof gives bit bounds, including the case of signed
coefficients. Its root-isolation import is
\citet[Theorem 36 in the linked preprint]{sagraloff2015}: integer univariate roots can be
isolated and refined in time polynomial in degree, coefficient length,
and requested precision bits. Rational denominator clearing and
square-free preprocessing have polynomial cost. Exact sign testing
then follows by evaluating at rational endpoints and between isolated
roots; values at roots themselves are zero. Gaussian positivity and
polynomial exactness are classical \citep[Section 3.5(v)]{dlmfquadrature}.

\begin{lemma}[Cumulative curvature and inverse quantiles]
\label{lem:certified-curvature}
Let $f$ be a densely encoded rational polynomial with
$f''\ge0$ nondecreasing on $[0,1]$, and let $\eta>0$ be rational.
For $F(x)=\int_0^x\rho_\eta(t)\,dt$, rational prefix queries $x$ and
requested absolute errors $\zeta>0$ admit certified rational estimates
in polynomial bit time. If $f$ is nonaffine, the normalized quantile
$F^{-1}(\theta F(1))$, for rational $\theta\in[0,1]$, admits a dyadic
approximation within any rational input tolerance $\delta>0$ in
polynomial bit time. The dependence on reciprocal tolerances is through
their bit lengths; the dependence on degree is numerical, as appropriate
for dense input. Endpoint quantiles are exact. All quantile outputs at
a fixed requested accuracy can use a common dyadic precision.
\end{lemma}
\begin{proof}
Put $H=f''/\eta$. Recognize $H=0$ exactly and return zero integrals;
the quantile assertion excludes this case. Otherwise set $d=\deg H$,
$d_0=\max(1,d)$, and $U=\max(1,H(1))$. Monotonicity and
nonnegativity imply $H(x)>0$ for $0<x\le1$, since a zero at a positive
point would force a polynomial to vanish on an interval. Both $\rho_\eta$
and $\sqrt H$ are bounded by $U$. We may replace $\zeta>1$ by one.
Choose $\lambda=2^{-J}$, $J\ge1$, with
$\lambda\le\zeta/(16U)$ and omit $[0,\lambda]$, losing mass at most
$\zeta/16$. The square-root branch is selected precisely where
\[
 P(x)=(1-x)^2H(x)-1\ge0.
\]
This is a nonzero polynomial because $P(1)=-1$. Isolate its distinct
roots in disjoint rational neighborhoods of total length at most
$\zeta/(16U)$. Omit those neighborhoods and determine the branch on
each remaining interval by a rational sign test. Tangencies and
multiple roots are handled by square-free isolation. The branch
$(1-x)H(x)$ has an exact rational antiderivative.

We first describe the simpler analytic panels available when $H$ has
nonnegative coefficients. Partition each dyadic interval
$[2^{-j-1},2^{-j}]$, $0\le j<J$, into $16d_0$ equal intervals.
For any subinterval of one panel, with center $c$ and length $\ell$,
$\ell/c\le1/(16d_0)$. On $|z-c|\le\ell$, every monomial
$(z/c)^k$ has argument at most $1/8$ in absolute value and modulus
less than two. Positive coefficients put $H(z)$ in an open right-half-plane
sector, and $|H(z)|\le2H(c)$. Thus its principal square root is
holomorphic there and bounded by $2U$. There are polynomially many
panels. This argument does not extend directly to signed coefficients.

For general coefficients, start with $[\lambda,1]$ and bisect a panel
$I=[a,b]$, center $c$, length $\ell$, until its exact rational Taylor
coefficients $H(c+t)=\sum_{k=0}^dh_k(c)t^k$ satisfy
\begin{equation}\label{eq:taylor-panel-certificate}
 \sum_{k=1}^d|h_k(c)|\ell^k\le H(c)/2.
\end{equation}
Accepted panels have $\Re H(z)\ge H(c)/2>0$ and
$|\sqrt{H(z)}|\le2U$ on the same disk. To bound the number of tests,
let $\alpha_1,\ldots,\alpha_d$ be the complex roots, with multiplicity;
they are used only in the proof. If every $|c-\alpha_j|\ge8d\ell$,
factorization gives
\[
 \frac{\sum_{k\ge1}|h_k(c)|\ell^k}{H(c)}
 \le\prod_j\left(1+\frac\ell{|c-\alpha_j|}\right)-1
 \le e^{1/8}-1<\tfrac12.
\]
Thus every failed center lies within $8d\ell$ of some root's real
part. Centers at one subdivision depth have spacing $\ell$, so at
most $d(16d+2)$ panels fail at that depth.

There are also only polynomially many depths. Clear denominators to
an integer polynomial $H_{\rm int}$ and set
$P_0=H_{\rm int}x(x-1)$, of degree $n=d+2$, with coefficients
bounded by $2^T$, $T\ge1$. Every root has modulus at most $2^{T+1}$.
The primitive square-free part has leading coefficient dividing that
of the primitive $P_0$, and all coefficients bounded by $2^{B_0}$,
where $B_0=(n+1)T+2n$ suffices by expansion in its roots. Its nonzero
integer discriminant gives the loose separation bound
\[
 |\alpha-\beta|\ge\sigma_*:=2^{-(B_0+2)n^2}
\]
for distinct roots: bound all other differences by $2^{T+2}$ and the
leading-coefficient factor by its height in the discriminant product.
A real root of $H$ is outside $(0,1]$; its distance to $[\lambda,1]$
is at least $\min(\lambda,\sigma_*)$. A nonreal root is separated
from its conjugate, so its imaginary part has modulus at least
$\sigma_*/2$. Consequently every root has distance at least
$\rho_*:=\min(\lambda,\sigma_*/2)$ from this interval.
At depth $K$ with $2^{-K}\le\rho_*/(8d_0)$ every panel passes.
Here $K$ is polynomial in the input and precision lengths. The bound
on failed panels per depth proves $O(d_0^2K+1)$ total tree nodes.
For constant $H>0$ the first panel passes. Exact Taylor tests at these
rational centers have polynomial coefficient length.

Intersect the accepted panels with the branch intervals and $[0,x]$.
If a subinterval has center $c'$ and length $\ell'$, then
$|c'-c|+\ell'\le(\ell+\ell')/2\le\ell$; its radius-$\ell'$ disk
remains in the certified disk. On such an interval an order-$q$
Gauss--Legendre rule has error at most $8U\ell'4^{-q}$.
Indeed the degree-$(2q-1)$ Taylor polynomial has uniform error at
most $4U4^{-q}$, and positivity, exactness, and weight sum $\ell'$
bound the integral and quadrature errors separately. Taking
$q=\lceil\frac12\log_2(128U/\zeta)\rceil$ makes their total at
most $\zeta/16$, since panel lengths sum to at most one.

We spell out rational conditioning of the rule. Let $P_q$ be the
Legendre polynomial on $[-1,1]$, with nodes $r_i$, and
$H_q=(q+1)(2q)!$. Its coefficient absolute sum is at most $H_q$.
The weights mapped to an interval of length one satisfy
\begin{equation}\label{eq:gauss-weight-conditioning}
 w_i=\frac1{(1-r_i^2)P_q'(r_i)^2},\qquad
 (qH_q)^{-2}\le w_i\le1.
\end{equation}
Polynomial exactness follows by dividing by $P_q$ and integrating
Lagrange cardinal polynomials; applying exactness to their squares
proves positivity, and the constant polynomial gives weight sum one.
The derivative bound $|P_q'|\le qH_q$ proves the lower weight bound.
The denominator in \eqref{eq:gauss-weight-conditioning} is at least
one. Polynomial derivative bounds and root refinement therefore give
positive rational weight estimates of relative error
$\zeta/(64U)$ using polynomially many bits, without constructing a
joint algebraic number field.

Put $V=\max(1,\sum_{k\ge1}k|H_k|)$, a rational derivative bound
on $[0,1]$. Rational evaluations of the monotone $H$ at the endpoints
of an isolated node interval enclose its value; their difference is at
most $V$ times the width. The inequality
$\sqrt v-\sqrt u\le\sqrt{v-u}$ for $0\le u\le v$ and rational
square-root bisection give absolute function-value error at most
$\zeta/64$ with polynomial precision. The positive weight sum bounds
all numerical quadrature errors by $\zeta/16$. Exact polynomial
integrals, rational sums, and denominator products retain polynomial
bit length. Together the cutoff, branch neighborhoods, analytic
quadrature, and rational evaluation cost less than $\zeta/4$, leaving
slack for a certified error-$\zeta$ answer.

Finally let $0<\delta\le1/2$; larger requests can be reduced to this
range. Set
\begin{equation}\label{eq:quantile-inverse-modulus}
 h_0=H(\delta/4)>0,\qquad
 \kappa=(\delta/4)\min(1,h_0),\qquad
 \mu=\delta\kappa/64.
\end{equation}
On $[\delta/4,1-\delta/4]$ both branches of $\rho_\eta$ are at
least $\kappa$. Every input interval of length $\delta$ contains a
subinterval of length $\delta/2$ in this range, so carries mass at
least $\delta\kappa/2$. The exact rational $h_0$ has polynomial
encoding even with cancellation. In the positive-coefficient case,
one can instead use $h_0=a(\delta/4)^k$ from any positive monomial
$a x^k$ of $H$, giving the same conclusion without cancellation.
Compute both $F(s)$ and $F(1)$ within $\mu$ in a dyadic bisection
for $F(s)-\theta F(1)=0$. The difference error is at most $2\mu$.
A certified nonzero sign updates the bracket. If the estimated
magnitude is at most $2\mu$, the true magnitude is at most
$4\mu<\delta\kappa/2$, which implies input distance less than
$\delta$ from the unique quantile. Otherwise
$\lceil\log_2(2/\delta)\rceil$ steps suffice. Endpoint quantiles
are handled exactly; padding early outputs gives common precision.
Since $\log(1/\mu)$ is polynomially bounded, so is the inverse
computation.
\end{proof}

\subsection{Mass-accurate knots and a seven-bit comparison}

For the formulation, accuracy in cumulative curvature mass suffices;
the knots need not approximate exact quantiles in input distance.
This permits a fixed absolute integration accuracy even where the
density is very small. A continuous polygonal path from zero to one
still covers the input interval if rounding reverses adjacent knots.

\begin{theorem}[Compiled monotone curvature]\label{thm:compiled-curvature}
For a densely encoded rational polynomial with nonnegative monotone
curvature on a nondegenerate rational interval and a positive rational
tolerance $\eta$, a polynomial-time algorithm constructs
a polynomial-size rational MILP graph relaxation of absolute error
$\eta$ using at most $\pconv(f;\eta)+7$ binaries.
For nondecreasing curvature on $[0,1]$, the actual count $L$ satisfies
\begin{equation}\label{eq:compiled-mass-count}
 2^L\le5M_\eta+2\le120N_\eta+2.
\end{equation}
Its error is at most $15\eta/16$. Affine functions use no binaries.
\end{theorem}
\begin{proof}
Normalize the rational interval and reflect its coordinate if curvature
is nonincreasing. Assume henceforth nonaffine nondecreasing curvature on
$[0,1]$. Put $\zeta=1/512$. Compute $\widehat M$ within $\zeta$
using \cref{lem:certified-curvature} and set
\[
 U=\max(\zeta,\widehat M+\zeta),\quad
 L=\max\{0,\lceil\log_2((5/2)U)\rceil\},\quad N=2^L.
\]
Then $M_\eta\le U\le M_\eta+1/256$ and $U/N\le2/5$.
The rational number $B=1+\sum_k k(k-1)|c_k|/\eta$ bounds the
density. For interior index $q$, bisect for mass target $t=qU/N$,
evaluating $F$ to error $\zeta$. Update the left bracket endpoint when
$\widehat F+\zeta<t$, and the right endpoint when
$\widehat F-\zeta>t$; stop
otherwise, with mass error at most $2\zeta$. After
$\lceil\log_2(B/\zeta)\rceil+1$ strict steps stop at the bracket
midpoint. When $t\le M_\eta$ the bracket contains a preimage and
has mass width at most $\zeta$. When $t>M_\eta$, every strict step
updates the left endpoint and its excess over total mass is at most $2\zeta$, so
the mass error is at most $3\zeta$. With forced $R_0=0,R_N=1$,
all outputs therefore satisfy the common safe certificate
\begin{equation}\label{eq:mass-knot-certificate}
 |F(R_q)-qU/N|\le4\zeta.
\end{equation}
Common dyadic precision is obtained by padding. No density lower bound
is used here, including for targets above the total mass.

Consecutive returned knots, whether ordered, reversed, or equal, enclose
mass at most $U/N+8\zeta\le133/320$. By
\eqref{eq:mass-local-remainder}, their chord error is at most
\[
 \eta\left[(133/320)^2+\tfrac32(133/320)\right]
 =\frac{81529}{102400}\eta<\frac{13}{16}\eta.
\]
Compute each rational endpoint value of $f$ downward to error at most
$\eta/8$. Dense rational evaluation and downward dyadic rounding have
polynomial bit cost. Add a fixed offset for negative values in the
compiler. If $y$ is their interpolant at the same continuous weight as
the input, then $-\eta/8\le y-f(x)\le13\eta/16$. Impose
\begin{equation}\label{eq:compiled-scalar-band}
 y-13\eta/16\le w\le y+\eta/8.
\end{equation}
Every true graph point is admitted: the polygonal input path begins
at zero and ends at one, hence covers the interval. Every admitted
point has error at most $15\eta/16$. Determinism gives identical
computed data for a shared endpoint. \Cref{lem:indexed-compiler}
uses exactly $L$ binaries, without enumerating the cells.

Since $M_\eta\le48\,2^{\pconv}$,
$(5/2)U\le120\,2^{\pconv}+5/512<128\,2^{\pconv}$, proving
the seven-bit count. Also $2^L\le\max(1,5U)\le5M_\eta+2$,
which proves \eqref{eq:compiled-mass-count}.
\end{proof}

For positive polynomials, endpoint values can alternatively be computed
monomial by monomial using the rounded binary-power procedure in
\cref{lem:binary-power-evaluation}. If $C=\sum_{k\ge2}c_k>0$,
individual downward error $\min(1,\eta/(8C))$ suffices; restore the
affine part exactly at the actual input. This implements the same band
and will be useful when numerical degrees are too large for exact
rational evaluation.

\subsection{Arbitrary convex dense polynomials}

\begin{theorem}[Hybrid convex-polynomial compiler]\label{thm:convex-hybrid}
Let a densely encoded rational polynomial $f$ be convex on $[0,1]$.
For rational $\eta>0$, a deterministic polynomial-time algorithm
constructs a polynomial-size rational MILP with error at most $\eta$
and at most $\pconv(f;\eta)+11$ binaries. Coefficients may have either
sign and curvature need not be monotone. Convexity can be a promise or
checked by exact univariate sign testing. Affine functions are exact
with zero binaries.
\end{theorem}
\begin{proof}
Assume degree $D\ge2$. Set
$V=\max(1,\sum_{k=1}^Dk|c_k|)\ge\max|f'|$ and choose
$h=2^{-B}$, $B\ge2$, with $h\le\eta/(32V)$. The implicit grid
has $2^B+1$ points, but its indices have polynomial length.
If a chord interval $[a,b]$ has error at most $\xi$, every subinterval
of $[\max(0,a-h),\min(1,b+h)]$ has error at most $\xi+2Vh$.
Indeed the old and expanded chords differ by at most $2Vh$ at $a,b$
and between them: the function is $V$-Lipschitz, endpoint value changes are at most
$Vh$, and every chord slope has absolute value at most $V$.
On each added strip the chord error is at most $2Vh$, and restriction
does not increase it. Round an optimal $\eta/4$ partition down to the
grid, keep zero and one, and discard repeated knots. Each resulting
interval lies within one of the original intervals expanded by $h$:
use the last original knot with its rounded left value. Thus there
is a grid partition with at most $N_{\eta/4}\le9N_\eta$ cells of
error at most $5\eta/16<\eta/2$.

Greedily extend each grid interval to the farthest feasible endpoint
at error $\eta/2$. Binary search uses $O(B)$ exact decisions of the
polynomial predicate $s_{a,b}-f-\eta/2\le0$ on $[a,b]$.
Root isolation and sign testing decide it in polynomial bit time,
including equality. Feasibility is monotone in the endpoint and a
one-step interval is always feasible. The greedy partition minimizes
count on this grid: by induction its $k$th endpoint is no earlier than
that of any feasible partition, since the remaining part of a comparison
interval is feasible by restriction. Its eventual count $G$ satisfies
\begin{equation}\label{eq:hybrid-greedy-count}
 G\le N_{\eta/4}\le9N_\eta.
\end{equation}
Run only until completion or $9D$ cells. If it finishes, enumerate
these polynomially many rational bands, using endpoint rounding and
\eqref{eq:compiled-scalar-band}. Here
$G\le27\,2^{\pconv}$, so at most $\pconv+5$ binaries suffice.

If it has not finished, \eqref{eq:hybrid-greedy-count} proves
$N_\eta>D$. Isolate the distinct roots of $f'''$ in $(0,1)$ in
disjoint rational brackets of width at most $h$, with nonroot rational
endpoints. Endpoint roots need no bracket; if $f'''=0$ none are needed.
Together the brackets and their complementary intervals form
$b\le3D$ nondegenerate pieces. Every bracket has chord error at most
$2Vh\le\eta/16$ and uses one cell. On each complement $f''$ is
monotone and nonnegative, so \cref{thm:compiled-curvature} applies
after rational affine normalization and, if needed, reflection.
All endpoints, sign decisions, and substitutions have polynomial bit
cost. The logarithm of root separation bounds the required precision;
no uniform grid at inverse root separation is enumerated.

Let $n_j$ be the optimal $\eta$ chord count on piece $j$. Splitting
a global optimal partition at its boundaries gives
$\sum_jn_j\le N_\eta+b-1$. Local cell counts are at most
$120n_j+2$, including the one-cell brackets. Therefore
\begin{equation}\label{eq:hybrid-total-cells}
 K\le120N_\eta+122b\le120N_\eta+366D
 <486N_\eta\le1458\,2^{\pconv}<2048\,2^{\pconv}.
\end{equation}
Compute local counts and cumulative sums in binary. A single global
index $0\le k<K$ selects a piece and its local index by comparisons
and subtraction; unused binary codes are excluded. There are only
polynomially many pieces, so running each indexed local routine and
selecting its output is polynomial time. Reversal maps local knots
back to their original rational interval. A common input denominator
is the product of piece-endpoint denominators times the largest local
dyadic denominator, of polynomial length. Endpoint values of $f$ are
rounded downward within $\eta/8$ using a fixed dyadic precision and
offset. Apply \cref{lem:indexed-compiler} and
\eqref{eq:compiled-scalar-band}. Each local path covers its piece;
the union covers $[0,1]$ with error at most $15\eta/16$.
The global index uses $\lceil\log_2K\rceil\le\pconv+11$ binaries.
\end{proof}

Greedy maximal segmentation is established in
\citet[Proposition 1 and Corollary 1]{codsi2025}. Their convex-corridor
algorithm already yields a continuous PWL approximation
(Remark 3), and computes each maximal segment with logarithmically
many function and derivative oracle calls at its requested endpoint
precision (Section 4.1). Their Lemma 4 bounds the extra segments caused
by splitting at changes of convexity by the number of split points.
Here the split instead makes the curvature monotone. The hybrid's
additional content is the lower-bound test that absorbs this splitting
cost, together with certified rational random access to the remaining
cells. It gives the eleven-bit comparison with arbitrary convex lifts
in polynomial bit time for dense input; continuity, greedy segmentation,
and logarithmic binary encoding are established ingredients.

\subsection{Separable sums and independent outputs}

\begin{lemma}[Jensen superadditivity]\label{lem:jensen-superadditivity}
For every continuous convex $\phi$ and ordered points
$a=x_0<\cdots<x_h=b$,
\begin{equation}\label{eq:jensen-superadditivity}
 J_\phi(a,b)\ge\sum_{j=0}^{h-1}J_\phi(x_j,x_{j+1}).
\end{equation}
For $\phi(x)=\sum_{k\ge2}c_kx^k$, $c_k\ge0$, its feature curve
$\psi(x)=(\sqrt{c_k}\,x^{k/2})_{k:c_k>0}$ also satisfies
\begin{equation}\label{eq:feature-jensen}
 \tfrac14\|\psi(a)-\psi(b)\|_2^2\le J_\phi(a,b)
 \le\tfrac12\|\psi(a)-\psi(b)\|_2^2,
\end{equation}
and its squared secant lengths are superadditive along ordered points.
\end{lemma}
\begin{proof}
For $C^2$ convex $\phi$,
$J_\phi(a,b)=\frac12\int_a^b\min(t-a,b-t)\phi''(t)\,dt$.
The long tent dominates each shorter tent, whose interiors are disjoint.
The same proof applies to a convex piecewise-linear function expressed
as an affine term plus nonnegative hinges. Uniform convex piecewise-linear
interpolation then proves \eqref{eq:jensen-superadditivity} for every
continuous convex function. For \eqref{eq:feature-jensen}, convexity
of $x^{k/2}$ gives the lower bound termwise, and
$((a+b)/2)^k\ge(ab)^{k/2}$ gives the upper. Feature increments have
nonnegative coordinates; expanding a squared sum of increments proves
the last statement. Irrational feature coefficients occur only in the
geometric proof.
\end{proof}

\begin{theorem}[Separable comparisons]\label{thm:separable-scalar}
On the product box $[0,1]^r$, consider either
$f(x)=\ell(x)+\sum_i\phi_i(x_i)$ with scalar absolute tolerance
$\eta$, or independent outputs $f_i(x)=\ell_i(x)+\phi_i(x_i)$
with separate tolerances $\eta_i$. Each $\ell_i$ is affine. For
continuous convex summands and unrestricted real-coefficient size,
\begin{equation}\label{eq:separable-finite-counts}
 \pbin\le\pconv+7r\quad\text{for sums},\qquad
 \pbin\le\pconv+4r\quad\text{for independent outputs}.
\end{equation}
For the computational statements, assume rational polynomial data,
including the affine terms, and positive rational tolerances.
For dense positive polynomials the polynomial rational
construction gives respectively $\pout\le\pconv+12r$ and
$\pout\le\pconv+9r$. For arbitrary rational dense convex polynomials,
it gives respectively $\pout\le\pconv+16r$ and
$\pout\le\pconv+13r$. Affine-only coordinates are retained exactly;
if $r=0$, no integers are needed.
\end{theorem}
\begin{proof}
At local tolerance $\tau_i$, take a maximal scalar incompatible set
of size $P_i$ as in \cref{lem:scalar-chords}. For independent outputs
use $\tau_i=\eta_i$. Any distinct product points violate some output's
midpoint tolerance, so $2^{\pconv}\ge\prod_iP_i$.
For sums take $\tau_i=\eta/r$. Adjacent selected points have gap
strictly greater than $\tau_i$; \cref{lem:jensen-superadditivity}
gives gap greater than $\tau_i h$ for index distance $h>0$.
Select product index vectors greedily, deleting the lattice ball of
$\ell_1$ radius $r$ after each selection. Such a ball has at most
\[
 \#\{z\in\Z^r:\|z\|_1\le r\}
 \le2^r\sum_{z\in\Z^r}2^{-\|z\|_1}=6^r
\]
points. The retained collection has at least $\prod_iP_i/6^r$
points with pairwise scalar gaps exceeding $\eta$. Thus
\begin{equation}\label{eq:separable-product-packing}
 2^{\pconv}\ge\prod_iP_i/6^r\quad\text{for sums}.
\end{equation}
These packings prove lower bounds and are not enumerated by the
construction.

Write $L_i$ for the actual local number of binaries. Finite chord
encoding has $2^{L_i}\le2N_{\tau_i}\le12P_i$. For the positive
curvature compiler,
$M_{\tau_i}\le2M_{2\tau_i}\le48N_{2\tau_i}\le96P_i$;
its definition gives $2^{L_i}\le\max(1,5U_i)\le481P_i$.
For the hybrid, either branch has at most $486N_{\tau_i}$ cells,
so $2^{L_i}\le972N_{\tau_i}\le5832P_i$.
Sum the logarithms. Independent outputs incur respectively
$\log_2 12<4$, $\log_2 481<9$, and $\log_2 5832<13$ per
coordinate. For sums \eqref{eq:separable-product-packing} adds
$\log_2 6$, yielding bounds below $7$, $12$, and $16$.
Compose the local graph relaxations and restore the affine outputs.
Local errors sum to at most $\eta$ in the sum model and remain within
their separate tolerances in the independent model. All dense
polynomial constructions have polynomial rational size and time;
replacing $\eta$ by $\eta/r$ adds only $O(\log r)$ precision bits.
\end{proof}

\section{Positive polynomial systems and error allocation}
\label{sec:positive-allocation}

Consider the multi-output system
\begin{equation}\label{eq:positive-polynomial-system}
 f_j(x)=\ell_j(x)+\sum_{i=1}^r\sum_{k=2}^{D_i}c_{jik}x_i^k,
 \qquad c_{jik}\ge0,\quad x\in[0,1]^r,
\end{equation}
where every coordinate is nonlinear in some output and $\ell_j$ is
affine. Inactive coordinates can be retained continuously. Let $K$ be
an unconditional error body, meaning it is invariant under every
coordinate sign change. Convexity then implies
\begin{equation}\label{eq:positive-unconditional-domination}
 u\in K,\quad |v_j|\le|u_j|\ (1\le j\le m)
 \quad\Longrightarrow\quad v\in K:
\end{equation}
the box generated by sign changes of $u$ lies in $K$.
Set $C_{ji}=\sum_kc_{jik}$ and
\begin{equation}\label{eq:positive-allocation}
 D_K=\max\left\{\prod_ip_i:0\le p_i\le1,\ Cp\in K\right\},
 \qquad \Phi_K=-\tfrac12\log_2D_K.
\end{equation}
There is a strictly positive feasible allocation because zero is interior
to $K$. An optimizer is positive. The case $r=0$ is affine and uses no
integers; all subsequent statements in this section assume $r\ge1$.
For a tolerance box the constraint is $\sum_iC_{ji}p_i\le\eta_j$.
For computational assertions data are rational, and $K$ has a polynomial
rational strong separation oracle with known positive rational inner
and outer radii. Its scalar specialization in
\cref{lem:block-logdet-oracle} returns an exactly feasible positive
rational $p$ with
\begin{equation}\label{eq:positive-allocation-oracle}
 \prod_ip_i\ge e^{-1}D_K.
\end{equation}
That proof includes rational logarithmic tangents, an explicit inner
ball, and the central-ball repair \eqref{eq:central-ball-repair}; it
requires neither nonnegative $C$ nor unconditional $K$. The stronger
hypotheses here are needed for graph-error geometry, not for allocation
optimization. Final formulations use explicit rational output bands,
not an exact linear representation of the potentially curved body $K$.

\subsection{A direct transformed-covariance and Taylor construction}

\begin{proposition}[Baseline polynomial allocation]\label{prop:positive-baseline}
Define
\[
 A_{r,D}=\log_2[\omega_r(r+2)^{r/2}]+\sum_i\log_2D_i-r/2,
 \qquad B_{r,D}=\sum_i\log_2D_i+r.
\]
Then \eqref{eq:positive-polynomial-system} satisfies
\begin{equation}\label{eq:positive-baseline-law}
 \max(0,\Phi_K-A_{r,D})\le\pconv\le\pbin
 \le\Phi_K+B_{r,D}.
\end{equation}
For rational dense input, the constructive count is at most
$\pconv+2\sum_i\log_2D_i+4r+1$, with polynomial rational size and
time. Its direct prefix construction is retained below because it
uses exact powers and an elementary Taylor rectangle.
\end{proposition}
\begin{proof}
For $2\le k\le D$ and $a,b\in[0,1]$, convexity of $x^{k/2}$ and
the $D/k$-Lipschitz property of $u^{D/k}$ give
\[
 J_{x^k}(a,b)\ge\tfrac14(a^{k/2}-b^{k/2})^2
 \ge D^{-2}(a^{D/2}-b^{D/2})^2.
\]
Map each compact parity support by the coordinate homeomorphism
$t_i=x_i^{D_i/2}$. These transformed supports still cover the unit
cube. On a positive-volume support, let $\Sigma$ be the covariance
of its uniform distribution and average over independent points.
The original Jensen vector belongs to $K$ and is nonnegative, and
its expectation componentwise dominates
$2\sum_iC_{ji}\Sigma_{ii}/D_i^2$. By convexity and
\eqref{eq:positive-unconditional-domination}, $p_i=2\Sigma_{ii}/D_i^2$ is
feasible; its caps follow from $\Sigma_{ii}\le1/4$.
The covariance-volume inequality \eqref{eq:covariance-volume} and
Hadamard's inequality give
\[
 \vol(S)\le\omega_r(r+2)^{r/2}\sqrt{\det\Sigma}
 \le\omega_r(r+2)^{r/2}\prod_i(D_i/\sqrt2)\sqrt{D_K}.
\]
A cover by at most $2^p$ such sets proves the lower bound. All volumes
are in the transformed cube; the homeomorphism is not asserted to
preserve volume. Zero-volume supports do not affect the argument.

For a positive allocation $p$, use original-axis depths
$L_i=\lceil\log_2(D_i/\sqrt{p_i})\rceil$ and widths
$h_i=2^{-L_i}$. Write $x_i=a_i+r_i$, where $a_i$ is the binary
prefix and $0\le r_i\le h_i$. Every output has the exact Taylor
expression
\[
 T_j=\ell_j(x)+\sum_{i,k}c_{jik}
                 (a_i^k+k a_i^{k-1}r_i).
\]
The coordinate second derivative is at most $D_i^2C_{ji}$, so
$0\le f_j-T_j\le(Cp)_j/2$. The rational rectangle
$T_j\le w_j\le T_j+(Cp)_j/2$ contains the graph and has
$|w-f(x)|\le Cp/2$, which implies error in $K$.

The prefix powers require no extra integers. With
$a=\sum_{\ell=1}^L2^{-\ell}z_\ell$, generate $v_0=1,t_0=r$ and
\begin{equation}\label{eq:exact-prefix-powers}
 v_k=\sum_\ell2^{-\ell}(z_\ell v_{k-1}),\qquad
 t_k=\sum_\ell2^{-\ell}(z_\ell t_{k-1}).
\end{equation}
Exact binary-times-bounded-continuous hulls, using $0\le v_k\le1$
and $0\le t_k\le h$, force $v_k=a^k,t_k=a^kr$. Generate through
$D$ and $D-1$, respectively, and share these variables across outputs.
This uses $O(D_iL_i)$ rows and continuous variables per coordinate.
The count is at most $-\frac12\log_2\prod_ip_i+B_{r,D}$.
An optimal allocation proves the finite upper. The rational allocation
\eqref{eq:positive-allocation-oracle} adds only $1/(2\ln2)$.
Since $A_{r,D}<\sum_i\log_2D_i+3r$ by the Gaussian volume bound,
the stated constructive comparison follows. All exact rational lengths
and recurrence counts are polynomial for dense input.
\end{proof}

Radix disaggregation and polynomial MILP approximation already appear in
\citet{teles2013}; the reused binary-product mechanism is established.
The substantive comparison here includes a lower bound against every
convex integer lift, not just a particular polynomial discretization.
For diagonal PSD quadratics the specialized sharper constants remain:
with $f_j=\frac12\sum_ih_{ji}x_i^2+\ell_j$, use $C_{ji}=h_{ji}$
in \eqref{eq:positive-allocation}. The expected Jensen vector is
$C\diag(\Sigma)/4$, and the shared square upper has
$|w-f|\le Cp/8$. Thus
\begin{equation}\label{eq:unconditional-diagonal-sharp}
 \max(0,\Phi_K-A_r')\le\pconv\le\pbin\le\Phi_K+r,
 \quad A_r'=r+\log_2[\omega_r(r+2)^{r/2}]<4r,
 \quad \pout\le\pconv+5r+1.
\end{equation}
This agrees with the structured quadratic analysis and extends its
coordinatewise bands to the whole unconditional body.

\subsection{Supporting scalarization removes degree from the lower bound}

\begin{theorem}[Double-logarithmic degree overhead]\label{thm:positive-loglog}
Put $J_i=\lceil\log_2D_i\rceil$,
$S_i=\lceil\log_2(J_i+1)\rceil$, and
$A_r=r/2+\log_2[\omega_r(r+2)^{r/2}]<7r/2$. Then
\begin{equation}\label{eq:positive-loglog-law}
 \max(0,\Phi_K-A_r)\le\pconv\le\pbin
 \le\Phi_K+r+\sum_iS_i.
\end{equation}
For rational dense input, a polynomial-time rational construction has
\begin{equation}\label{eq:positive-loglog-count}
 \pout\le\pconv+\tfrac92r+\sum_iS_i+1.
\end{equation}
The same constructive bound holds for sparse nonnegative polynomial
input with exponents in binary, with size and time polynomial in the
sparse encoding, by the endpoint implementation below.
\end{theorem}
\begin{proof}
Let $p^*$ maximize $\sum_i\log p_i$. Strict feasibility of a small
positive common allocation and the convex normal-cone sum rule give
an outward normal $\lambda\in N_K(Cp^*)$ and $\mu_i\ge0$, where
$N_K(z)=\{v:v^T(y-z)\le0\text{ for all }y\in K\}$, such that
\begin{equation}\label{eq:allocation-supporting-normal}
 1/p_i^*=(C^T\lambda)_i+\mu_i,\qquad
 \mu_i(p_i^*-1)=0.
\end{equation}
We may choose $\lambda\ge0$. At a positive coordinate of $Cp^*$,
downward closure permits a small decrease, forcing the corresponding
normal component nonnegative. A zero coordinate is a zero row of $C$,
since $C\ge0$ and $p^*>0$. Set its normal component to zero: the
support function of an unconditional body is unchanged by sign changes
and monotone in absolute coordinates, so the normal inequality remains
valid while its value at $Cp^*$ is unchanged.
Set $a_i=(C^T\lambda)_i\ge0$ and
$b=\lambda^TCp^*=h_K(\lambda)$. The scalar allocation
$0\le p\le1$, $\sum_i a_ip_i\le b$ has the same optimum $D_K$.
Indeed, concavity and \eqref{eq:allocation-supporting-normal} give
\[
 \sum_i\log(p_i/p_i^*)
 \le\sum_i(p_i-p_i^*)/p_i^*
 =a^T(p-p^*)+\sum_i\mu_i(p_i-p_i^*)\le0.
\]
Points with zero product cannot improve the optimum. If $\lambda=0$,
all $p_i^*=1$ and the lower bound is trivial. An index with $a_i=0$
also has $p_i^*=1$.

For $a_i>0$, normalize
$\phi_i(x)=a_i^{-1}\sum_k(\sum_j\lambda_jc_{jik})x^k$ and let
$g_i=\sqrt{\phi_i}$. This is an increasing homeomorphism of $[0,1]$
onto itself and is convex: it is the Euclidean norm of nonnegative
convex coordinates $\sqrt{d_{ik}}x^{k/2}$ with $\sum_kd_{ik}=1$.
For $a_i=0$, use $g_i(x)=x$. Since a nonnegative convex $g$ satisfies
$J_{g^2}(x,y)\ge(g(x)-g(y))^2/4$, each parity support transformed
by $t_i=g_i(x_i)$ obeys
$\frac14\sum_i a_i(t_i-s_i)^2\le b$. Taking independent uniform
points gives $\frac12\sum_i a_i\Sigma_{ii}\le b$.
Thus $p_i=\Sigma_{ii}/2$ is feasible for the scalar allocation and
$\prod_i\Sigma_{ii}\le2^rD_K$. The covariance-volume bound now
gives $\vol(S)\le2^{r/2}\omega_r(r+2)^{r/2}\sqrt{D_K}$.
Coverage of the transformed cube proves the degree-independent lower
bound. Multipliers and transformed coordinates are proof devices;
the construction need not compute or rationalize them.

For the upper, partition each coordinate into $J_i+1$ layers
\begin{equation}\label{eq:dyadic-endpoint-layers}
 [1-2^{-j},1-2^{-j-1}]\ (0\le j<J_i),\qquad
 [1-2^{-J_i},1].
\end{equation}
If a layer has left endpoint $\ell$ and width $s$, then for every
$2\le k\le D_i$ its normalized monomial curvature is
\begin{equation}\label{eq:layer-curvature}
 k(k-1)s^2x^{k-2}\le2.
\end{equation}
On the final layer $s\le1/D_i$ gives a bound of one. Else $s\le1/2$
and $x\le1-s$. For $k=2$ the bound is at most $1/2$; for $k\ge3$
put $v=(k-2)s$, giving
$k(k-1)s^2(1-s)^{k-2}\le(v+1)^2e^{-v}\le4/e<2$.

Encode each layer using $S_i$ binaries and continuous selectors
$\theta_j\ge0$ summing to one, each bounded by the matching bit
literals. At integral bits the matching selector is one and all others
zero; unused codes are infeasible. Such implied-binary selectors can
be used in exact binary-product hulls without additional declared
integers. For allocation $p_i$, set
$L_i=\lceil\frac12\log_2(1/p_i)\rceil$, $h_i=2^{-L_i}$.
Write the normalized layer input as $a_0+u$, with $L_i$ prefix bits
and $0\le u\le h_i$. In original coordinates impose
\[
 a=\sum_j\theta_j(\ell_j+s_ja_0),\quad
 v=\sum_js_j\theta_ju,\quad x=a+v.
\]
For any bounded variable $q$, the product $aq$ is represented as
$\sum_j[\ell_j(\theta_jq)+s_j(\theta_j(a_0q))]$;
$a_0q$ uses the prefix bits and selector products use their forced
binary values. Generate exact $a^k,a^kv$ by
\eqref{eq:exact-prefix-powers} with this implementation. Bounds
$0\le a^k\le1$, $0\le a^kv\le h_i$ suffice. At $L_i=0$, the
prefix is zero and those products disappear. The size per coordinate
is $O(D_i(L_i+J_i+1)+J_iS_i)$, polynomial for dense input.

The resulting exact Taylor expression $T_j$ satisfies
$0\le f_j-T_j\le\sum_iC_{ji}h_i^2\le(Cp)_j$ by
\eqref{eq:layer-curvature}. Use $T_j\le w_j\le T_j+(Cp)_j$.
The graph is contained and every error belongs to $K$ by
\eqref{eq:positive-unconditional-domination}. Counting the $L_i+S_i$ bits
proves the finite upper in \eqref{eq:positive-loglog-law}.
The allocation \eqref{eq:positive-allocation-oracle} adds at most
$1/(2\ln2)$ to the count. Combining with $A_r<7r/2$ proves
\eqref{eq:positive-loglog-count}. The sparse implementation is proved
next and uses exactly the same layer and cell bits.
\end{proof}

\begin{lemma}[Rounded binary powers]\label{lem:binary-power-evaluation}
For dyadic $x\in[0,1]$, integer $k\ge1$ in binary, and rational
$\xi>0$, downward-rounded repeated squaring returns $A\in[0,1]$
with $0\le x^k-A\le\xi$ in polynomial time in the input length,
$\log k$, and precision encoding. It suffices to use $P$ fractional
bits at least the input precision and $\lceil\log_2(k/\xi)\rceil$.
\end{lemma}
\begin{proof}
All rounded factors stay below their exact values in $[0,1]$.
If factors representing exponents $a,b\ge1$ have errors at most
$(a-1)2^{-P}$ and $(b-1)2^{-P}$, multiplication and one downward
rounding have error at most $(a+b-1)2^{-P}$. The exact base has zero
error. Induction over $O(\log k)$ products proves error at most
$(k-1)2^{-P}$. An initial exact accumulator one requires no rounding
when multiplied by an already representable factor. All arithmetic
uses polynomially many bits.
\end{proof}

\begin{proof}[Sparse construction for \cref{thm:positive-loglog}]
Let $E_i$ list the exponents, in binary. Obtain the same rational
allocation, layer count $J_i$, and depth $L_i$. The only integer inputs
are the $S_i$-bit layer index $j\le J_i$ and $L_i$-bit cell index $q$.
Compute the exact dyadic endpoints
$a=\ell_j+s_jqh_i$, $b=\ell_j+s_j(q+1)h_i$ at common precision
$J_i+L_i+1$. For each listed $k$, compute downward endpoint powers
$A_{ik},B_{ik}$ with error at most $\xi_i=p_i/4$ using
\cref{lem:binary-power-evaluation}; a common precision
$P_i\ge\max(J_i+L_i+1,\lceil\log_2(D_i/\xi_i)\rceil)$ suffices.
Only listed exponents are evaluated. Compile all endpoint values and
interpolate with one common weight per input, giving
$x_i=(1-\theta_i)a+\theta_i b$ and
$z_{ik}=(1-\theta_i)A_{ik}+\theta_i B_{ik}$.
The exact input grid covers the coordinate interval. The curvature
bound \eqref{eq:layer-curvature} and the chord bound yield
\[
 -\xi_i\le z_{ik}-x_i^k\le h_i^2/4.
\]
Set $T_j=\ell_j(x)+\sum_{i,k}c_{jik}z_{ik}$ and impose
\begin{equation}\label{eq:sparse-polynomial-band}
 T_j-\sum_iC_{ji}h_i^2/4\le w_j
 \le T_j+\sum_iC_{ji}\xi_i.
\end{equation}
These bands contain the graph and have
$|w-f|\le Cp/2$, hence error in $K$. All circuit computations are
polynomial in sparse exponent lengths and allocation length; no exact
numerator of a numerical-degree power is formed. By
\cref{lem:indexed-compiler} the count remains $\sum_i(L_i+S_i)$.
\end{proof}

\section{Pure powers and inverse interpolation}\label{sec:pure-powers}

For one power per coordinate shared across all outputs,
\begin{equation}\label{eq:pure-power-system}
 f_j(x)=\ell_j(x)+\sum_{i=1}^r c_{ji}x_i^{\alpha_i},
 \qquad c_{ji}\ge0,
\end{equation}
the degree penalty disappears. We retain both an explicit rational
reciprocal construction for dense integer degrees and an indexed
construction polynomial in binary rational exponent lengths. The
continuous descriptions differ even where their integer counts agree.
Throughout, $K$ is unconditional and each active coordinate has some
positive coefficient; $r=0$ is affine and exact.

\begin{lemma}[Power-coordinate chords]\label{lem:power-chords}
For $\alpha\ge2$, a chord of $x^\alpha$ on $[a,b]\subset[0,1]$
has error at most $4(b^{\alpha/2}-a^{\alpha/2})^2$.
Moreover $J_{x^\alpha}(a,b)\ge(a^{\alpha/2}-b^{\alpha/2})^2/4$.
For all $\alpha>1$, set $s=\min(1,\alpha-1)$. Then
\begin{equation}\label{eq:scaled-power-jensen}
 J_{x^\alpha}(a,b)\ge\tfrac{s}{8}
                         (b^{\alpha/2}-a^{\alpha/2})^2.
\end{equation}
On a uniform power-coordinate cell $[t,t+h]$, $t=kh$, interpolating
$(t^{2/\alpha},t^2)$ and $((t+h)^{2/\alpha},(t+h)^2)$ gives
\begin{equation}\label{eq:scaled-power-chord}
 0\le y-x_0^\alpha\le4s h^2.
\end{equation}
\end{lemma}
\begin{proof}
The stronger Jensen bound for $\alpha\ge2$ follows from convexity
of $x^{\alpha/2}$. Write $A=a^{\alpha/2}$, $B=b^{\alpha/2}$.
If $A\le B/2$, the chord is at most $B^2\le4(B-A)^2$.
Otherwise the usual curvature chord bound and the mean-value estimate
give error at most
\[
 \frac{\alpha(\alpha-1)b^{\alpha-2}}8(b-a)^2
 \le\frac{\alpha-1}{2\alpha}(b/a)^{\alpha-2}(B-A)^2
 \le2(B-A)^2,
\]
since $(b/a)^{\alpha/2}<2$. This proves the first assertion and
\eqref{eq:scaled-power-chord} for $\alpha\ge2$.

If $1<\alpha<2$ and $b>0$, the curvature on $(a,b)$ is at least
$\alpha(\alpha-1)b^{\alpha-2}$. Subtract its quadratic and apply
convexity, extending to $a=0$ by continuity. Since
$u^{\alpha/2}\ge u$ on $[0,1]$,
$b^{\alpha/2}-a^{\alpha/2}\le b^{\alpha/2-1}(b-a)$.
This proves \eqref{eq:scaled-power-jensen}; the zero interval is
immediate. For the first power-coordinate cell,
$y-x_0^\alpha=h^2(\theta-\theta^\alpha)
\le s h^2\theta(-\log\theta)\le s h^2/e$.
For another cell put $u=t+h\le2t$ and $\beta=2/\alpha\in(1,2)$.
Its input length is at most $\beta u^{\beta-1}h$ and its maximum
curvature is $\alpha s t^{2-2\beta}$. The chord error is at most
$[s/(2\alpha)](u/t)^{4/\alpha-2}h^2\le2s h^2$, proving the
safe bound \eqref{eq:scaled-power-chord}.
\end{proof}

\begin{proposition}[Finite pure-power dimension]\label{prop:pure-power-finite}
For real exponents $\alpha_i\ge2$, with $C=(c_{ji})$ and the
allocation \eqref{eq:positive-allocation},
\begin{equation}\label{eq:pure-power-finite}
 \max(0,\Phi_K-A_r)\le\pconv\le\pbin\le\Phi_K+2r,
 \qquad A_r=r/2+\log_2[\omega_r(r+2)^{r/2}]<7r/2.
\end{equation}
Thus $\pbin\le\pconv+11r/2$ with unrestricted real coefficients
and size. For all $\alpha_i>1$, set $s_i=\min(1,\alpha_i-1)$,
$C^{\rm eff}_{ji}=c_{ji}s_i$, and define $\Phi_{\rm eff}$ by
\eqref{eq:positive-allocation} with $C^{\rm eff}$ in place of $C$.
Then
\begin{equation}\label{eq:scaled-power-lower}
 \pconv\ge\Phi_{\rm eff}-A_r',
 \qquad A_r'=r+\log_2[\omega_r(r+2)^{r/2}]<4r.
\end{equation}
\end{proposition}
\begin{proof}
Map parity supports by $t_i=x_i^{\alpha_i/2}$, an increasing
homeomorphism for every positive exponent. For $\alpha_i\ge2$, the
nonnegative Jensen vector dominates $C[(t-u)^2]/4$. Independent
uniform points in a transformed support of covariance $\Sigma$ give
$C\diag(\Sigma)/2\in K$. Thus $p_i=\Sigma_{ii}/2$ is feasible,
and the covariance-volume proof gives the first lower bound. The
scaled Jensen inequality instead gives
$C^{\rm eff}\diag(\Sigma)/4\in K$; using
$p_i=\Sigma_{ii}/4$ proves \eqref{eq:scaled-power-lower}.
The caps follow from $\Sigma_{ii}\le1/4$ in both cases.

For the finite upper, use a positive allocation and power-coordinate
widths $h_i=2^{-L_i}$ with
$L_i=\lceil\frac12\log_2(1/p_i)\rceil+1$.
The knots $(kh_i)^{2/\alpha_i}$ give chord error at most
$4h_i^2\le p_i$ by \cref{lem:power-chords}. Encode each union of
chord bands with $L_i$ binaries and set the outputs to the positive
combination of these coordinate approximations. Then
$|w-f|\le Cp\in K$ in the domination sense, and the exact graph is
contained. Sum the counts at an optimal allocation. The generally
irrational knots and potentially exponential row count do not supply
a compact rational construction.
\end{proof}

For rational dense integer degrees $D_i$, the direct original-axis
prefix construction also retains the bound
$\pout\le\Phi_K+\sum_i\log_2D_i+r+1/(2\ln2)$, hence
$\pout\le\pconv+\sum_i\log_2D_i+9r/2+1$ by the stronger
pure-power lower bound. The inverse constructions below remove this
degree contribution while changing the continuous representation.

\subsection{Positive Stieltjes approximation with rational coefficients}

Positive resolvent approximations of fractional powers are established.
For example, substituting $\lambda=1/t$ in
\citet[Section 3.3, equation (37) and Lemma 3.4]{bonito2015} gives
positive sums of $t/(t+b)$ approximating $t^\gamma$. The next proof
supplies explicit rational coefficient lengths, exact endpoint
normalization, and uniform parameter bounds for the present construction.

\begin{lemma}[Certified rational inverse power]\label{lem:stieltjes-power}
Given an integer $D\ge3$ and rational $0<\delta<1/2$, one can
construct positive rational $a_k,b_k$ such that
\[
 R(t)=\sum_{k=1}^M a_k\frac{t}{t+b_k},\qquad
 R(0)=0,\quad R(1)=1,\quad
 \sup_{0\le t\le1}|R(t)-t^{2/D}|\le\delta.
\]
The term count is
$M=O(D(1+\log(1/\delta)+\log D)^2)$; time and total coefficient
length are polynomial in $D$ and the precision encoding. For $D=2$
use $R(t)=t$ exactly.
\end{lemma}
\begin{proof}
Put $\gamma=2/D\in(0,1)$ and use the Stieltjes integral
\[
 I_\gamma(t)=\int_0^\infty\frac{t}{t+s}s^{\gamma-1}\,ds,
 \qquad I_\gamma(0)=0.
\]
Substitution gives $I_\gamma(t)=t^\gamma I_\gamma(1)$, and splitting
at one gives $2\le I_\gamma(1)\le D/2+3$.
Let $e=\min(\delta,1/4)$, let $p$ be the least integer with
$2^{-p}\le e$, and put $d_*=\lceil\log_2D\rceil$,
\[
 L=\lceil(D/2)(p+d_*+8)\rceil,\quad U=3(p+8),\quad N=L+U.
\]
The omitted tails below $2^{-L}$ and above $2^U$ are bounded,
uniformly for $0\le t\le1$, by
$2^{-L\gamma}/\gamma\le e/512$ and
$2^{-U(1-\gamma)}/(1-\gamma)\le3e/256$, respectively.
Their sum is less than $e/64$.

On each dyadic panel substitute $s=2^jv$, $1\le v\le2$,
$-L\le j<U$. The integrand is
$2^{j\gamma}v^{\gamma-1}t/(t+2^jv)$.
On $|v-3/2|\le1$ it is holomorphic and has modulus at most four:
$\Re v\ge1/2$, $|v^{\gamma-1}|\le2$; for $j\le0$ the other
factors have product modulus at most one, and for $j\ge0$ use
$|t/(t+2^jv)|\le2^{1-j}$. The $t=0$ integrand is zero.
The degree-$(2q-1)$ Taylor remainder on $[1,2]$ is at most
$8\,4^{-q}$. Gaussian positivity and exactness give panel error
at most $16\,4^{-q}$. Taking
$q=\lceil[p+\lceil\log_2(256N)\rceil]/2\rceil$
bounds the sum of these errors by $e/16$.
Writing the common nodes and weights as $v_i\in(1,2)$, $w_i>0$,
$\sum_iw_i=1$, yields the positive sum
\[
 Q(t)=\sum_{j=-L}^{U-1}\sum_{i=1}^q A_{ji}\frac{t}{t+B_{ji}},
 \quad A_{ji}=2^{j\gamma}w_iv_i^{\gamma-1},\quad B_{ji}=2^jv_i.
\]
It satisfies $\|Q-I_\gamma\|_\infty<5e/64$ and
$0\le Q(t)\le Q(1)\le D+4$.

Choose $\tau=e/[128(D+4)]$. Approximate every $A_{ji},B_{ji}$
by positive rationals with relative error at most $\tau$. A positive
term changes relatively by at most $4\tau$, because its perturbed
ratio lies between $(1-\tau)/(1+\tau)$ and its reciprocal.
Consequently $\|\widehat Q-Q\|_\infty\le e/32$ and
$E:=\|\widehat Q-I_\gamma\|_\infty<7e/64<e/8$.

We verify polynomial precision, including small weights. With
$H_q=(q+1)(2q)!$, \eqref{eq:gauss-weight-conditioning} gives
\[
 2^{-L}\le B_{ji}\le2^U,\qquad
 \frac{2^{-L}}{2(qH_q)^2}\le A_{ji}\le2^U.
\]
All logarithmic magnitude bounds are polynomial in $L,U,q$.
Root isolation and rational interval evaluation of the Legendre weight
formula therefore give the required relative precision in polynomial
bit time. For the remaining powers use
$A_{ji}=(w_i/v_i)(2^jv_i)^{2/D}$: certified bisection of a positive
$D$th root solves $z^D=(2^jv_i)^2$. All arguments have explicit
positive upper and lower bounds of polynomial bit length. Derivatives
of this fixed-depth expression and the root map on those bounds have
logarithmic magnitudes polynomial in $D,L,U,q$; refining its input
intervals and squared comparisons thus costs polynomially many bits.
Raising rational test values to degree $D$ has polynomial time and
length in numerical $D$, as asserted. No unbounded-precision
real-arithmetic operation is assumed.

Normalize exactly by the positive rational $\widehat Q(1)$:
$a_{ji}=\widehat A_{ji}/\widehat Q(1)$, $b_{ji}=\widehat B_{ji}$.
The endpoints of $R=\widehat Q/\widehat Q(1)$ are exactly zero
and one, and
\[
 |R(t)-t^\gamma|\le\frac{2E}{2-E}\le e\le\delta.
\]
Positive coefficients make $R$ strictly increasing. Rational sums
and normalization retain polynomial total encoding, since denominator
products add their bit lengths. Finally $N=O(D(p+\log D))$ and
$q=O(p+\log D)$, giving the stated term count. This degree dependence
is compatible with dense input, not with a claim polynomial in $\log D$.
\end{proof}

\begin{lemma}[Rational endpoint products with existing bits]
\label{lem:rational-endpoint-products}
Let $a=\sum_{\ell=1}^L2^{-\ell}z_\ell$, $h=2^{-L}$, and
$R=P/Q$ for supplied rational polynomials with a supplied certificate
$Q(t)\ge q_{\min}>0$ on $[0,1]$. For either $t=a$ or $t=a+h$
and $0\le\theta\le1$, the value $\theta R(t)$ has an exact linear
representation using only the $L$ existing binaries and
$O((d+1)(L+1))$ rows and continuous variables, where
$d=\max(\deg P,\deg Q)$, with the zero numerator treated directly.
The coefficient length is polynomial in the supplied data, including
$q_{\min}$. Signs of polynomial coefficients are unrestricted.
\end{lemma}
\begin{proof}
For $t=a+sh$, $s\in\{0,1\}$ fixed, impose
\[
 0\le v_k\le1/q_{\min},\quad
 v_k=\sum_\ell2^{-\ell}(z_\ell v_{k-1})+shv_{k-1}
 \ (1\le k\le d),\quad
 \sum_{k=0}^dQ_kv_k=\theta.
\]
Products use exact bounded binary-product hulls. Recurrences force
$v_k=t^kv_0$, so positivity of $Q(t)$ forces
$v_0=\theta/Q(t)$. These values satisfy the bounds; therefore
$\sum_kP_kv_k=\theta R(t)$ exactly. The same proof covers $L=0$,
where $a=0,h=1$ and the prefix products are absent. Using weights
$1-\theta,\theta$ represents the interpolation of $R(a),R(a+h)$.
If $R(0)=0,R(1)=1$, its polygonal path covers $[0,1]$ even when
$R$ is nonmonotone; retain $0\le x\le1$ to clip any excursions.
A uniform error $\|R-g\|_\infty\le\delta$ implies error at most
$\delta$ between interpolated endpoint values. The lemma supplies
neither a denominator-certification algorithm nor a short rational
approximation for an arbitrary $g$.
\end{proof}

\begin{theorem}[Dense reciprocal construction]\label{thm:reciprocal-powers}
For \eqref{eq:pure-power-system} with integer degrees $D_i\ge2$,
rational dense input and the allocation oracle assumptions, a rational
MILP of polynomial size is constructible in polynomial time with
$\pout\le\pconv+13r/2+1$.
\end{theorem}
\begin{proof}
For a positive feasible rational allocation, set
$L=\lceil\frac12\log_2(1/p)\rceil+2$, $h=2^{-L}$,
$\delta=p/(16D)$ in each coordinate. Take the rational approximation
of \cref{lem:stieltjes-power}. Each term has reciprocal variables
\[
 (a+b_k)v_k^- =1-\theta,\qquad
 (a+h+b_k)v_k^+=\theta,\qquad 0\le v_k^\pm\le1/b_k.
\]
Expanding $a$ in its bits makes these exact linear constraints via
bounded binary products. They impose
\[
 x=\sum_k a_k[1-b_k(v_k^-+v_k^+)]
   =(1-\theta)R(a)+\theta R(a+h).
\]
For $D=2$ use $x=a+h\theta$ instead. The auxiliary
$y=a^2+2ha\theta+h^2\theta$ is linearized with the same bits,
since $a^2=\sum_\ell2^{-\ell}z_\ell a$ and similarly for
$a\theta$. The size is $O((M+1)(L+1))$, and reciprocal bounds
have polynomial coefficient length. No new integers occur.

Let $x_0=(1-\theta)a^{2/D}+\theta(a+h)^{2/D}$.
Then $|x-x_0|\le\delta$ and \cref{lem:power-chords} gives
$0\le y-x_0^D\le4h^2$. Since $x^D$ is $D$-Lipschitz,
$-D\delta\le y-x^D\le4h^2+D\delta$. Impose the rational band
\begin{equation}\label{eq:inverse-power-band}
 y-(4h^2+D\delta)\le q\le y+D\delta.
\end{equation}
The increasing endpoint-normalized $R$ covers every input, so the
whole graph is admitted, while $|q-x^D|\le4h^2+2D\delta\le3p/8$.
Set $w_j=\ell_j(x)+\sum_ic_{ji}q_i$; unconditionality controls
all errors. Using \eqref{eq:positive-allocation-oracle}, the count
is at most $\Phi_K+3r+1/(2\ln2)$. Combine with
\cref{prop:pure-power-finite}. Polynomial numerical-degree complexity
follows from the positive rational approximation lemma.
\end{proof}

\subsection{Rational powers with binary exponent encoding}

\begin{lemma}[Random access to inverse powers]\label{lem:inverse-power-index}
For an $L$-bit dyadic grid and a rational exponent $\alpha>1$ in
binary, one can compute dyadic knots $R_k\in[0,1]$ within rational
error $0<\delta<1/2$ of $(k/2^L)^{2/\alpha}$ in time polynomial
in $L$, exponent encoding, and precision encoding. The endpoints are
exact, shared indices give identical outputs, and all outputs may
use one fixed precision. For integer exponents there is also a direct
rounded-power bisection algorithm.
\end{lemma}
\begin{proof}
The case $L=0$ consists only of the two exact endpoints. For the direct
integer procedure let $D\ge2$, $0<k<2^L$, $u=(k/2^L)^2\ge2^{-2L}$,
and $\tau=\delta2^{-2L}/16$. Bisection for $s^D=u$ uses the enclosure
$[A,A+\tau]$ from \cref{lem:binary-power-evaluation}, at precision
$P\ge\max(R+1,\lceil\log_2(D/\tau)\rceil)$, where
$R=\lceil\log_2(2/\delta)\rceil$. A separated enclosure gives a
strict sign. If it contains $u$, then $|s^D-u|\le\tau$ and both
power values are at least $u/2$. On the intervening segment,
$Dz^{D-1}\ge z^D\ge u/2$, so the mean-value theorem gives
$|s-u^{1/D}|\le2\tau/u\le\delta/8$. Otherwise $R$ steps make
the bracket sufficiently short. Padding outputs to $R+1$ fractional
bits works. This requires $O(\log(1/\delta))$ bisections,
$O(\log D)$ rounded products per step, and
$P=O(\log D+L+\log(1/\delta))$ bits, even for enormous $D$.

For rational $\alpha$, put $\beta=2/\alpha\in(0,2)$ and normalize
the interior dyadic $t$ as $v=2^et\in[1,2)$, $1\le e\le L$.
For $w\in[1,2]$, define
\[
 H_N(w)=2\sum_{j=0}^{N-1}\frac{((w-1)/(w+1))^{2j+1}}{2j+1}.
\]
Integrating the geometric series gives
$0\le\log w-H_N(w)\le3\,9^{-N}$.
Choose $N$ with $6(L+1)9^{-N}\le\delta/8$ and compute the rational
$A=\beta(H_N(v)-eH_N(2))$. Then $-2L\le A\le0$ and
$|A-\beta\log t|\le\delta/8$. Let $M$ be a power of two with
$4L\le M<8L$, so $q=A/M\in[-1/2,0]$. For an odd $J$ large
enough that $2^{-(J+1)}\le\delta/(16M)$, the exact rational Taylor
sum $P_J(q)=\sum_{j=0}^Jq^j/j!$ obeys
\[
 \tfrac12\le P_J(q)\le e^q\le1,
 \qquad 0\le e^q-P_J(q)\le\delta/(16M).
\]
The first lower bound follows by pairing the terms after $1+q$;
the other bounds are the alternating-series estimate.
The $M$th-power map is $M$-Lipschitz on $[0,1]$, and the exponential
is one-Lipschitz on the nonpositive axis. Compute $P_J(q)^M$ exactly
and round downward with error at most $\delta/4$. Its total error
from $t^\beta$ is at most $\delta/16+\delta/8+\delta/4<\delta$.
All series have $O(\log(L+1)+\log(1/\delta))$ terms. Exact rational
powering enlarges their bit lengths by at most $M=O(L)$, so the
algorithm has polynomial bit cost. For $\alpha\ge2$, one may use
the smaller constants $3(L+1)$, $-L\le A\le0$, and $2L\le M<4L$.
Endpoints and common output precision are handled separately as above.
\end{proof}

\begin{theorem}[Binary rational-exponent construction]\label{thm:rational-powers}
For rational \eqref{eq:pure-power-system} with binary-encoded
$\alpha_i>1$ and the unconditional allocation oracle assumptions,
a polynomial-time construction gives a rational MILP of polynomial
total encoding length with
\begin{equation}\label{eq:rational-power-count}
 \pout\le\pconv+7r+1.
\end{equation}
If all $\alpha_i\ge2$, the stronger count
$\pout\le\pconv+13r/2+1$ holds. Size depends polynomially on the
exponent bit lengths, not their numerical magnitudes.
\end{theorem}
\begin{proof}
For general $\alpha_i>1$ allocate with $C^{\rm eff}_{ji}=c_{ji}s_i$,
$s_i=\min(1,\alpha_i-1)$. Set
$L_i=\lceil\frac12\log_2(1/p_i)\rceil+2$,
$h_i=2^{-L_i}$, and $\delta_i=s_ip_i/(16\alpha_i)$.
The indexed algorithm computes both inverse knots and the exact
power-coordinate squares $(kh_i)^2,((k+1)h_i)^2$. Apply
\cref{lem:indexed-compiler} to impose their interpolants $(x_i,y_i)$
with one weight. Relative to the exact inverse interpolant $x_{0i}$,
$|x_i-x_{0i}|\le\delta_i$. By
\eqref{eq:scaled-power-chord} and the $\alpha_i$-Lipschitz bound,
\[
 -\alpha_i\delta_i\le y_i-x_i^{\alpha_i}
 \le4s_ih_i^2+\alpha_i\delta_i.
\]
The band
$y_i-(4s_ih_i^2+\alpha_i\delta_i)\le q_i
\le y_i+\alpha_i\delta_i$ therefore has error at most
$3s_ip_i/8$. Its graph containment follows from the continuous
polygonal path joining exact endpoint knots zero and one; reversed
or duplicate interior segments cause no problem. Assemble outputs
as positive combinations of the $q_i$. Every error is dominated by
$3C^{\rm eff}p/8$, and so lies in $K$.
The count is at most $\Phi_{\rm eff}+3r+1/(2\ln2)$; use
\eqref{eq:scaled-power-lower}. If all exponents are at least two,
take $s_i=1$ and the stronger lower bound in
\cref{prop:pure-power-finite}. All precisions, including
$\log(1/s_i)$ when $\alpha_i$ approaches one, have polynomial
rational encoding. Only the external cell bits are integral.
\end{proof}

\section{Relative error and rational encoding boundaries}
\label{sec:encoding-boundaries}

The preceding counts concern positive absolute error. Two distinct
limitations arise if the error vanishes at a domain endpoint or if
short rational coefficients are required for a sharply varying root.

\begin{proposition}[Relative accuracy near zero]\label{prop:relative-power}
For $q>1$ and any finite $\eta\ge0$, no convex lift with finitely
many integer coordinates can contain all $(x,x^q)$, $0<x\le1$,
and require $w\le(1+\eta)x^q$ at every admitted point. Thus finite
integer dimension is impossible for two-sided pure relative error.
For $0<q<1$, two-sided relative error is impossible with finite integer
dimension if $0\le\eta<1$, whereas $\pconv=0$ if $\eta\ge1$.
Including $x=0$ does not change these conclusions.
For fixed $q>1$, $\eta>0$ on the truncated domain $[a,1]$,
\begin{equation}\label{eq:relative-truncated-law}
 \pconv(a),\ \pbin(a)=\Theta_{q,\eta}(\log\log(1/a))
 \qquad(a\downarrow0).
\end{equation}
\end{proposition}
\begin{proof}
For $q>1$ choose an integer $M\ge2$ with
$M^{q-1}>2^q(1+\eta)$. In a purported $p$-integer lift choose
exact witnesses for $M^p+1$ points of the sequence $M^{-k}$.
Two witness indices have the same residue modulo $M$. Label their
inputs $0<x<y$, so $x/y\le1/M$. The convex combination of weight
$1/M$ on the witness at $y$ has an integer index and visible coordinates
\[
 u=y/M+(1-1/M)x\le2y/M,\qquad
 v=y^q/M+(1-1/M)x^q\ge y^q/M.
\]
Thus $v/u^q\ge M^{q-1}/2^q>1+\eta$, a contradiction. The proof
uses only finite convex combinations and applies to nonclosed lifts
and unbounded integer witnesses.

For $0<q<1$, $\eta<1$, choose $M$ with
$M^{q-1}+M^{-q}<1-\eta$ and use the sequence $M^{-2k}$.
The same residue argument gives $x/y\le M^{-2}$ and $u\ge y/M$,
so $v/u^q\le M^{q-1}+M^{-q}<1-\eta$, again a contradiction.
If $\eta\ge1$, the convex set
$\{(x,w):0<x\le1,\ 0\le w\le x^q\}$ contains the graph and
has relative error at most one. Adding $(0,0)$ handles the closed domain.

On $[a,1]$, the $N+1$ inputs $1,M^{-1},\ldots,M^{-N}$,
$N=\lfloor\log_M(1/a)\rfloor$, must have distinct residues.
Therefore $\pconv(a)\ge\log_M(N+1)$. For the upper let
$\rho=(1+\eta)^{1/q}>1$ and divide $[a,1]$ into at most
$K=\max(1,\lceil\log(1/a)/\log\rho\rceil)$ geometric intervals
$[\ell,u]$ with $u/\ell\le\rho$. Each rectangle
$\ell\le x\le u$, $\ell^q\le w\le u^q$ contains the graph
and has $1/(1+\eta)\le w/x^q\le1+\eta$, which implies the
required two-sided error. Encode the finite union with
$\lceil\log_2K\rceil$ binaries. These bounds prove
\eqref{eq:relative-truncated-law}; no matching leading constants
are asserted. The real-coefficient upper uses $O(\log(1/a))$ cells.
At zero error on any nondegenerate interval, strict convexity makes
arbitrarily many distinct graph points pairwise midpoint-incompatible,
so finite integer dimension is again impossible.
\end{proof}

The residue proof is an elementary extension of the midpoint method
\citep{lubin2022}. Affine slices transfer its convex-power obstruction
to nonnegative bilinear and multilinear products by setting all inputs
equal, and to quadratic-over-linear graphs on a fixed positive-denominator
slice approaching zero numerator. A positive absolute error floor or
truncation away from zero removes this specific obstruction.

\subsection{A fixed-error root graph requires long rational MILPs}

\begin{theorem}[Tight rational encoding barrier]\label{thm:root-encoding}
For integer $D\ge2$, let $f_D(x)=x^{1/D}$ on $[0,1]$.
Every rational MILP whose projection contains its graph and has
absolute vertical error at most $1/16$ has ordinary explicit total
binary encoding length $\Omega(D)$. This includes arbitrary numbers
of continuous auxiliaries and unbounded general integer variables.
Conversely a rational MILP with four binaries and total encoding
$O(D)$ achieves this error. Thus minimum rational MILP encoding is
$\Theta(D)$, although four binaries suffice for every $D$.
\end{theorem}
\begin{proof}
Let $s$ be total explicit encoding, counting rows, variables, and
rational coefficients. Freeze any integer witness $z_*$ for the
exact graph point $(2^{-D},1/2)$. Add $w\ge1/4$ and minimize $x$
over the resulting rational polyhedron. This LP is feasible, has
finite attained optimum $v$, and
\begin{equation}\label{eq:root-small-lp-value}
 0<v\le2^{-D}.
\end{equation}
Indeed the tube bounds inputs below by zero, while an attained point
with $x=0,w\ge1/4$ would violate its error bound at zero.

Clear denominators separately in each original row and the new row.
If $\ell_i$ is the sum of coefficient numerator and denominator bit
lengths in row $i$, including integer columns and right-hand side,
and $t_i$ its nonzero count, then $\sum_i\ell_i=O(s)$ and
$\sum_it_i=O(s)$. Continuous coefficients after clearing are integers
of magnitude at most $2^{\ell_i}$. Fixing $z_*$ changes only the
integer right-hand side, whose numerator can be arbitrarily large.
Split free continuous variables into nonnegative differences and add
slacks. The resulting integer row norms are at most
$\sqrt{2t_i+1}\,2^{\ell_i}$. After dependent equations are removed,
the feasible standard-form LP has an optimal basic solution; this
argument does not require a vertex in the original free-variable
polyhedron. Cramer's rule makes its positive objective value have
a denominator dividing the nonzero integer basis determinant.
Hadamard bounds every such determinant by
\[
 \log_2|\det B|\le\sum_i\ell_i+\tfrac12\sum_i\log_2(2t_i+1)
 \le C_0s
\]
for a universal encoding constant $C_0$. Thus $v\ge2^{-C_0s}$,
independently of witness magnitude. With
\eqref{eq:root-small-lp-value}, this proves $s\ge D/C_0$.
The row-wise total-length accounting is essential to this linear
bound in $D$.

For the upper choose rational knots $a_j=(j/16)^D$, $0\le j\le16$.
On segment $j$ impose
\[
 x=(1-\theta)a_j+\theta a_{j+1},\quad
 t=(j+\theta)/16,\quad t\le w\le t+1/16,
 \qquad0\le\theta\le1.
\]
Concavity and monotonicity give $t\le f_D(x)\le t+1/16$.
The segments cover $[0,1]$. Their finite union uses four binaries;
explicitly, continuous selectors $\lambda_j$ bounded by their four
matching literals sum to one, and $0\le v_j\le\lambda_j$ gives
\[
 x=\sum_j[a_j\lambda_j+(a_{j+1}-a_j)v_j],\qquad
 t=\sum_j(j\lambda_j+v_j)/16.
\]
The row and variable counts are constant, and each rational knot has
$O(D)$ bits. This proves the upper encoding bound.
\end{proof}

This is a denominator bound for rational linear systems, not a lower
bound on real-coefficient formulation size. The binary exponent $1/D$
has $O(\log D)$ bits, whereas the point of height one half is at
input $2^{-D}$. A small integer count alone cannot ensure a short
rational graph formulation.

\subsection{The same four bits with a short conic formulation}

A mixed-integer second-order cone formulation (MISOCP) uses rational
linear equations and second-order cone constraints with integer
variables. Its count is compared below with rational MILP encoding;
it is not included in the definition of $\pbin$.

\begin{lemma}[A fixed conic optimum as a coefficient]\label{lem:conic-value}
Suppose the primal-dual conic pair
\[
 \min\{c^Tz:Az=b,\ z\in\mathcal K\},\qquad
 \max\{b^Tu:A^Tu+s=c,\ s\in\mathcal K^*\}
\]
has attained equal finite optimum $v$. Then, for $\theta\ge0$,
the conic linear constraints
\begin{equation}\label{eq:conic-value-homogeneous}
 Az=\theta b,\ z\in\mathcal K,\qquad
 A^Tu+s=\theta c,\ s\in\mathcal K^*,\qquad
 c^Tz=b^Tu=t
\end{equation}
project exactly to $t=\theta v$. They introduce no bilinear constraints.
\end{lemma}
\begin{proof}
For $\theta>0$, divide the variables by $\theta$ and use weak
duality; equality of the objective values forces $t=\theta v$.
Scale attained optimizers for the converse. At $\theta=0$, use any
original feasible $z_0,(u_0,s_0)$ to get
$c^Tz=s_0^Tz\ge0$ and $b^Tu=-z_0^Ts\le0$. Their equality forces
$t=0$, which is feasible using zero variables. Thus no nonzero
recession value enters at zero weight.
\end{proof}

\begin{theorem}[Rational MILP--MISOCP encoding separation]
\label{thm:root-conic-separation}
For $D=2^B$, $B\ge1$, the fixed-error graph relaxation in
\cref{thm:root-encoding} has a rational MISOCP with exactly four
binaries and $O(B)$ constant-size numerical coefficients in a structured
cone-list description. Ordinary sparse matrix encoding has
$O(B\log(B+2))$ total bits. Every rational MILP for the same graph
error has $\Omega(2^B)$ total encoding.
\end{theorem}
\begin{proof}
For fixed rational $0<a<1$, the SOCP
\[
 \min z_B,\qquad z_0=a,\qquad z_i\ge z_{i-1}^2\quad(1\le i\le B)
\]
has optimum $v_B=a^{2^B}$, attained by $v_i=a^{2^i}$.
Each squaring constraint is
$(z_i,1/2,z_{i-1})\in\mathcal Q_r$, where
$\mathcal Q_r=\{(u,v,w):u,v\ge0,\ 2uv\ge w^2\}$ is the
self-dual rotated second-order cone. The rational point
$z_i=a+i(1-a)/(B+1)$ is strictly feasible. There is also an explicit
dual certificate: put $\lambda_B=1$ and
$\lambda_i=2v_i\lambda_{i+1}$ for $1\le i<B$.
The cone vectors
$\lambda_i(1,2v_{i-1}^2,-2v_{i-1})$ lie in $\mathcal Q_r$.
Their pairings with $(z_i,1/2,z_{i-1})$ sum to $z_B-v_B$;
coefficients of intermediate $z_i$ cancel, and the constants telescope
using $\lambda_iv_i=2\lambda_{i+1}v_{i+1}$.
They are nonnegative and vanish at the primal optimizer, certifying
attained equal dual value. These possibly long certificate coordinates
are witnesses, not coefficients in the SOCP. Conversion to conic
standard form has linear size. Apply \cref{lem:conic-value} to encode
$t=\theta a^{2^B}$ with $O(B)$ short rational conic data. Values
$a=0,1$ use direct equations.

Use the sixteen-segment construction of \cref{thm:root-encoding},
with four index bits and selectors $\lambda_j$. Write
$0\le v_j\le\lambda_j$, $r_j=\lambda_j-v_j$, $s_j=v_j$.
Two copies of the conic-value gadget impose
$X_j=(j/16)^D r_j$ and $Y_j=((j+1)/16)^D s_j$.
Set
\[
 x=\sum_j(X_j+Y_j),\qquad
 t=\sum_j(j\lambda_j+v_j)/16,\qquad t\le w\le t+1/16.
\]
Only the selected segment contributes; its graph containment and error
were proved above. There are constantly many $O(B)$ gadgets, and
all numerical data have constant length. Sparse variable and row
indices contribute at most $O(B\log(B+2))$ bits. The MILP lower
bound is \cref{thm:root-encoding} with $D=2^B$.
\end{proof}

Conic strong duality, repeated-squaring lifts, and short conic systems
with long solution encodings are established phenomena
\citep{boyd2004,odonnell2017,wang2024}. The separation above concerns
fixed vertical graph accuracy with the same four declared binaries.
It uses exact cone feasibility and exact primal-dual value equality;
it asserts no tolerance-stable numerical solution or polynomial-time
exact-output algorithm. In particular, a polynomial-size rational MILP
outer approximation of these conic constraints that still contains the
whole exact graph must violate the required vertical error somewhere,
by \cref{thm:root-encoding}.
